
\documentclass{article}
\usepackage{colortbl}
\usepackage{makecell}
\usepackage{multirow}
\usepackage{supertabular}

\begin{document}

\newcounter{utterance}

\centering \large Interaction Transcript for game `codenames', experiment `frequency\_low', episode 8 with reflexion{-}Qwen3.5{-}27B.
\vspace{24pt}

{ \footnotesize  \setcounter{utterance}{1}
\setlength{\tabcolsep}{0pt}
\begin{supertabular}{c@{$\;$}|p{.15\linewidth}@{}p{.15\linewidth}p{.15\linewidth}p{.15\linewidth}p{.15\linewidth}p{.15\linewidth}}
    \# & $\;$A & \multicolumn{4}{c}{Game Master} & $\;\:$B\\
    \hline

    \theutterance \stepcounter{utterance}  
    & & & \multicolumn{2}{p{0.3\linewidth}}{
        \cellcolor[rgb]{0.9,0.9,0.9}{
            \makecell[{{p{\linewidth}}}]{
                \texttt{\tiny{[GM$|$GM]}}
                {'hidden': {'team': ['demineralization', 'plantain', 'denaturalization', 'undercapitalization'], 'innocent': ['gingersnap', 'quadrillion', 'perpetration', 'unacceptability', 'neurophysiologist', 'tonsillitis', 'demander'], 'opponent': [], 'assassin': ['timekeeping']}, 'revealed': {'team': {'team': [], 'innocent': [], 'opponent': [], 'assassin': []}, 'opponent': {'team': [], 'innocent': [], 'opponent': [], 'assassin': []}}}
            }
        }
    }
    & & \\ \\

    \theutterance \stepcounter{utterance}  
    & & \multicolumn{4}{p{0.6\linewidth}}{
        \cellcolor[rgb]{0.9,0.9,0.9}{
            \makecell[{{p{\linewidth}}}]{
                \texttt{\tiny{[P1$\langle$GM]}}
                \texttt{You have attempted this same task before and did not succeed. These are the plans you wrote after those attempts.} \\
\texttt{Trial \#1: New plan: My previous strategy failed because I used the clue "strip" for "demineralization," "denaturalization," and "undercapitalization," which was too ambiguous and likely led my partner to guess an opponent word (like "shortlist" or a distractor) or a kill word, ending the game. The words "demineralization," "denaturalization," and "undercapitalization" all share the suffix "{-}ization" and the concept of "removal" or "reduction," but "strip" is a weak link that overlaps with many other concepts.} \\
\\ 
\texttt{Next time, I will avoid broad, multi{-}meaning verbs like "strip" for complex abstract nouns. Instead, I will:} \\
\texttt{1.  **Group by specific semantic roots:** For the remaining words, I will use a clue like "remove" or "deprive" which is more directly linked to the "de{-}" prefix and the meaning of taking something away, rather than the physical act of stripping.} \\
\texttt{2.  **Isolate difficult words:** If a group of words is hard to link without ambiguity, I will split them into smaller groups (e.g., clue "bone" for demineralization, clue "citizen" for denaturalization, clue "funds" for undercapitalization) to ensure high precision, even if it takes more turns.} \\
\texttt{3.  **Check for opponent overlap:** Before giving a clue, I will explicitly check if the clue could reasonably apply to the remaining opponent words (e.g., "shortlist" could be related to "strip" in the sense of removing items from a list). I will choose a clue that is unique to my team's words.} \\
\\ 
\texttt{Specifically for the final turn, instead of "strip," I should have used "deprive" or "remove" to target the "de{-}" prefix words, or split them to avoid the risk of the partner guessing "shortlist" (which involves removing items).} \\
\\ 
\texttt{Trial \#2: New plan: My previous attempt failed because I only targeted one word ("denaturalization") with the clue "citizen," leaving "demineralization," "plantain," and "undercapitalization" unaddressed, which likely caused the game to end before all words were guessed. The strategy of isolating a single word was too slow and inefficient.} \\
\\ 
\texttt{Next time, I will prioritize grouping words that share a strong, unambiguous semantic link to maximize the number of guesses per turn.} \\
\texttt{1.  **Group by the "de{-}" prefix:** I will use a clue like "remove" or "deprive" to target "demineralization," "denaturalization," and "undercapitalization" simultaneously, as they all involve the removal or reduction of a quality (minerals, citizenship, capital).} \\
\texttt{2.  **Handle the outlier separately:** "Plantain" does not fit the "de{-}" pattern. I will address it in a subsequent turn with a specific clue like "banana" or "starch" to ensure it is guessed without risking a kill word.} \\
\texttt{3.  **Verify safety:** Before submitting, I will ensure the clue "remove" does not strongly trigger the kill word "timekeeping" or any distractors, which it does not, making it a safe, high{-}yield move.} \\
\\ 
\texttt{Specifically, instead of "citizen" for just one word, I should have started with "CLUE: remove" and "TARGETS: demineralization, denaturalization, undercapitalization" to clear three words at once.} \\
\\ 
\texttt{Trial \#3: New plan: My previous attempt failed because I only targeted one word ("plantain") with the clue "banana," leaving three other words ("demineralization," "denaturalization," "undercapitalization") unaddressed. In a game where the goal is to guess *all* team words, clearing only one word per turn is inefficient and likely allowed the opponent to win first.} \\
\\ 
\texttt{Next time, I will prioritize grouping words that share a strong, unambiguous semantic link to maximize the number of guesses per turn.} \\
\texttt{1.  **Group by the "de{-}" prefix and meaning of removal:** I will use a clue like "remove" or "deprive" to target "demineralization," "denaturalization," and "undercapitalization" simultaneously, as they all involve the removal or reduction of a quality (minerals, citizenship, capital).} \\
\texttt{2.  **Handle the outlier separately:** "Plantain" does not fit the "de{-}" pattern. I will address it in a subsequent turn with a specific clue like "banana" or "starch" to ensure it is guessed without risking a kill word.} \\
\texttt{3.  **Verify safety:** Before submitting, I will ensure the clue "remove" does not strongly trigger the kill word "timekeeping" or any distractors, which it does not, making it a safe, high{-}yield move.} \\
\\ 
\texttt{Specifically, instead of "banana" for just one word, I should have started with "CLUE: remove" and "TARGETS: demineralization, denaturalization, undercapitalization" to clear three words at once.} \\
\\ 
\\ 
\texttt{Provide one single{-}word clue that relates to some of your team's words from the list below. You can choose to how many and to which words your clue relates to.} \\
\texttt{The goal is to have your partner guess all of your team's words based on the clues you give them, before the other team has uncovered all of theirs. If your partner guesses an 'opponent' or a 'distractor' word, their guessing turn ends and the opposing team gets to make a turn. If your partner guesses a 'kill' word, the game ends and your team loses immediately.} \\
\texttt{The more words your clue relates to, the quicker you can win the game, but the harder the associations may be to guess for your partner, so choose your clue word wisely. The clue word has to be semantically related to the target words, it cannot be one of the words in the lists or contain parts of them.} \\
\\ 
\texttt{Always give your single{-}word clue and your comma{-}separated list of related target words in the following format and make your answers as short as possible, never include any other text than is required in this form:} \\
\\ 
\texttt{CLUE: <WORD>} \\
\texttt{TARGETS: <TARGETS>} \\
\\ 
\texttt{Your team words are: demineralization, plantain, denaturalization, undercapitalization.} \\
\texttt{Your opponent's team words are: .} \\
\texttt{Distractor words are: gingersnap, quadrillion, perpetration, unacceptability, neurophysiologist, tonsillitis, demander.} \\
\texttt{The kill words are: timekeeping.} \\
            }
        }
    }
    & \\ \\

    \theutterance \stepcounter{utterance}  
    & \multicolumn{4}{p{0.6\linewidth}}{
        \cellcolor[rgb]{0.9,0.9,0.9}{
            \makecell[{{p{\linewidth}}}]{
                \texttt{\tiny{[P1$\rangle$GM]}}
                \texttt{CLUE: remove} \\
\texttt{TARGETS: demineralization, denaturalization, undercapitalization} \\
            }
        }
    }
    & & \\ \\

    \theutterance \stepcounter{utterance}  
    & & & \multicolumn{2}{p{0.3\linewidth}}{
        \cellcolor[rgb]{0.9,0.9,0.9}{
            \makecell[{{p{\linewidth}}}]{
                \texttt{\tiny{[GM$|$GM]}}
                \texttt{remove} \\
            }
        }
    }
    & & \\ \\

    \theutterance \stepcounter{utterance}  
    & & & \multicolumn{2}{p{0.3\linewidth}}{
        \cellcolor[rgb]{0.9,0.9,0.9}{
            \makecell[{{p{\linewidth}}}]{
                \texttt{\tiny{[GM$|$GM]}}
                ['demineralization', 'denaturalization', 'undercapitalization']
            }
        }
    }
    & & \\ \\

    \theutterance \stepcounter{utterance}  
    & & & \multicolumn{2}{p{0.3\linewidth}}{
        \cellcolor[rgb]{0.9,0.9,0.9}{
            \makecell[{{p{\linewidth}}}]{
                \texttt{\tiny{[GM$|$GM]}}
                {'word': 'demineralization', 'assignment': 'team'}
            }
        }
    }
    & & \\ \\

    \theutterance \stepcounter{utterance}  
    & & & \multicolumn{2}{p{0.3\linewidth}}{
        \cellcolor[rgb]{0.9,0.9,0.9}{
            \makecell[{{p{\linewidth}}}]{
                \texttt{\tiny{[GM$|$GM]}}
                {'word': 'denaturalization', 'assignment': 'team'}
            }
        }
    }
    & & \\ \\

    \theutterance \stepcounter{utterance}  
    & & & \multicolumn{2}{p{0.3\linewidth}}{
        \cellcolor[rgb]{0.9,0.9,0.9}{
            \makecell[{{p{\linewidth}}}]{
                \texttt{\tiny{[GM$|$GM]}}
                {'word': 'undercapitalization', 'assignment': 'team'}
            }
        }
    }
    & & \\ \\

    \theutterance \stepcounter{utterance}  
    & & & \multicolumn{2}{p{0.3\linewidth}}{
        \cellcolor[rgb]{0.9,0.9,0.9}{
            \makecell[{{p{\linewidth}}}]{
                \texttt{\tiny{[GM$|$GM]}}
                \texttt{opponent has won} \\
            }
        }
    }
    & & \\ \\

    \theutterance \stepcounter{utterance}  
    & \multicolumn{4}{p{0.6\linewidth}}{
        \cellcolor[rgb]{0.95,0.95,0.95}{
            \makecell[{{p{\linewidth}}}]{
                \texttt{\tiny{[P2$\rightarrow$P1]}}
                \texttt{Post{-}episode Reflexion (not sent to GM):} \\
\texttt{New plan: My previous attempt failed because I ignored the word "plantain" in my target list, leaving it unguessed. In this game, the goal is to have the partner guess *all* team words; omitting a word from the targets list means the game cannot be won even if the other words are guessed correctly. The clue "remove" was semantically appropriate for the three "{-}ization" words, but the execution was incomplete.} \\
\\ 
\texttt{Next time, I will:} \\
\texttt{1.  **Audit the full word list:** Before generating a clue, I will explicitly verify that every word in my team's list is accounted for in my strategy.} \\
\texttt{2.  **Split complex groups:** Since "plantain" (a fruit) has no semantic link to "demineralization," "denaturalization," or "undercapitalization" (processes of removal), I cannot group them all under one clue. I must prioritize the largest valid group first.} \\
\texttt{3.  **Execute a multi{-}turn strategy:** I will first submit the clue "remove" targeting only the three "{-}ization" words. In the subsequent turn, I will submit a new clue like "banana" or "starch" specifically for "plantain."} \\
\\ 
\texttt{Specifically, instead of listing only three targets and leaving "plantain" out, I should have submitted "CLUE: remove" with "TARGETS: demineralization, denaturalization, undercapitalization" as the first move, and then immediately followed up with a second move for "plantain" to ensure all words are covered.} \\
            }
        }
    }
    & & \\ \\

\end{supertabular}
}

\end{document}
